\subsection{The original degenerate orbital operator}
\label{shared-original-degenerate-operator}

The flag fibre estimate now supplies the local operator bound used
for the nonidentity terms with zero Gram determinant.  We retain
the original Haar measures, quotient measures, compact
representations and complex phases.  The argument below proves
the exact local assertion of the pinned
source~\cite{global-gram-degenerate-operator-source}.

Use the notation and normal-form hypotheses of
Theorem~\ref{thm:original-degenerate-fibre-estimate}.  Thus
$1\le j\le r$, $j+2\le k$, $D=|D_j(\gamma)|>0$,
$\Delta\ge1$ bounds the inverses of the nonzero leading minors,
and the common-stabilizer measure is the prescribed transport of
$s\,d\rho$, with $s>0$.  Put $a=k-1-j\ge1$.

For a negative $z\in W$, let $H_z$ be its stabilizer in
$U(J_W)$, embedded as the identity on $E_0\oplus F_0$.
It is compact; give it probability Haar measure $dh_z$.
Let $\pi:H_z\to U(\mathcal V)$ and
$\pi':H_{z'}\to U(\mathcal V')$ be continuous irreducible
representations on nonzero finite-dimensional complex Hilbert
spaces.  Suppose continuous functions $\Phi:G\to\mathcal V$ and
$\Phi':G\to\mathcal V'$ satisfy
\[
 \Phi(hnx)=\psi(n)\pi(h)\Phi(x),\qquad
 \Phi'(h'nx)=\psi'(n)\pi'(h')\Phi'(x),
 \qquad |\psi(n)|=|\psi'(n)|=1.
\]
Here $h\in H_z$, $h'\in H_{z'}$, $n\in N$, and $x\in G$.

Fix the original Haar measure $dg$ of $G$.  Use the original
quotient measures $\nu_z,\nu_{z'},\nu_R$, with the Weil formulas
\[
 \int_G f(g)\,dg
   =\kappa\int_G\int_{H_z}\int_N f(hnx)\,dn\,dh_z\,d\nu_z(x)
   =\int_G\int_{R_\gamma}f(\rho x)\,d\mu_R(\rho)\,d\nu_R(x)
\]
for nonnegative measurable $f$, and the corresponding formula
for $z'$.  The same $\kappa>0$ is used on both sides.  This
formulation permits the original realization of a quotient measure
as a measure on $G$.
Assume, for one $c_\Phi>0$,
\[
 \int_G\|\Phi\|^2d\nu_z=c_\Phi\dim\mathcal V,\qquad
 \int_G\|\Phi'\|^2d\nu_{z'}=c_\Phi\dim\mathcal V'.
\]
For an integer $q\ge16r^3+4r+2$ and $H\ge1$, assume the
original length-tail bounds
\begin{align*}
 \int_{\{S_z\ge tH\}}\|\Phi\|^2d\nu_z
       &\le(4r+4)t^{-2q}c_\Phi\dim\mathcal V,\\
 \int_{\{S_{z'}\ge tH\}}\|\Phi'\|^2d\nu_{z'}
       &\le(4r+4)t^{-2q}c_\Phi\dim\mathcal V'
                                     \qquad(t>0).
\end{align*}
Theorem~\ref{global-gram-original-all-length-tail} supplies
these bounds with the smaller coefficient $2r+2$ in the actual
arithmetic family.

Put
\[
 \sigma=\max\left(1,\frac{|D_l(\gamma)|}{(2H)^{2r}}:
                                               0\le l\le r\right).
\]
For $u\in\mathcal V$ and $v\in\mathcal V'$ write
$[u,v]w=\langle w,v\rangle u$, with the inner product linear
in its first argument.

\begin{theorem}[The actual local degenerate bound]
\label{thm:actual-local-degenerate-bound}
Under these hypotheses the original rank-one kernel is
$\nu_R$-integrable, and the original operator
\[
             J_\gamma=\frac1{c_\Phi}
                \int_G[\Phi(\gamma x),\Phi'(x)]\,d\nu_R(x)
\]
satisfies
\begin{align}
 \|J_\gamma\|_{\rm op}\le{}&
 \frac{\kappa}{s}\,
 16(4r+4)(4r)^{4r^2}(2H)^{16r^3+4r}
                       \Delta^{8r^2+2}\notag\\
    &{}\cdot\left(\frac{8\pi\mathrm e}{a}\right)^{ja}
    D^{-a}\,
    \sigma^{-(q-16r^3-4r-1)/r}.
 \label{eq:actual-local-degenerate-bound}
\end{align}
The assertion keeps the given quotient realization and every
complex phase.  Its norm estimate uses the absolute scalar
pairings, which are invariant under $R_\gamma$.
\end{theorem}

\begin{proof}
We supply the measure-theoretic and geometric inputs of the
already proved dyadic estimate (OD3)~\cite{global-gram-degenerate-dyadic}.

\emph{The measure when the two sides are interchanged.}
Conjugation by $\gamma$ sends $R_\gamma$ to $R_{\gamma^{-1}}$.
It sends $\mu_R$ to the original normal-form measure for
$\gamma^{-1}$ with the same scalar $s$.  To check the scalar,
first work at $g_0$.  On the first $j$ coordinates its block
matrix, in $E_j\oplus F_j$, is
\[
                  \begin{pmatrix}0&Z^{-*}\\ Z&0\end{pmatrix}.
\]
Within an antidiagonal block, conjugation sends an upper
unitriangular $u$ to $(Zu^{-1}Z^{-1})^*$.  Inversion in
upper unitriangular coordinates has absolute real Jacobian one.
The remaining operations are a permutation, conjugation, and
entrywise multiplications by ratios of entries of $Z$.
Those ratios have absolute value one within a block.
On the last block the action on $u$ is the identity.
If $m_0=g_0|_W$, the top radical coordinates transform as
\[
                    B_{\rm top}\longmapsto
                         \overline c\,m_0B_{\rm top},
                  \qquad S_{\rm top}\longmapsto S_{\rm top}.
\]
Since $m_0\in U(J_W)$ and $|c|=1$, their real Jacobian is
one.  Hence conjugation by $g_0$ preserves $d\rho$.
The prescribed transports by $n_1,n_2$ now prove the assertion
for $\gamma=n_1g_0n_2$.
The inverse cross matrix is $Z_\gamma^*$, so all leading-minor
absolute values, including $D$, are unchanged.  The quotient
measure on the reversed side is exactly
$(x\mapsto\gamma x)_*\nu_R$, by left invariance of $dg$
and the Weil formula.

\emph{A cutoff for the given quotient measure.}
There is a nonnegative continuous $\chi$ on $G$ such that
\[
               \int_{R_\gamma}\chi(\rho x)\,d\mu_R(\rho)=1
                                              \qquad(x\in G).
\]
One construction also verifies the convergence needed here.
Choose compactly supported nonnegative continuous $f_i$ whose
positive sets cover $G$.  Their averages
$F_i(x)=\int_{R_\gamma}f_i(\rho x)\,d\mu_R(\rho)$ are finite
and continuous, and bounded on compact sets: for $x$ in a
fixed compact set, the possible $\rho$ lie in a fixed compact
subset of the closed group $R_\gamma$.
Choose positive coefficients tending to zero fast enough that
both $\sum_i\epsilon_i f_i$ and $\sum_i\epsilon_i F_i$
converge uniformly on each member of a compact exhaustion.
Their sums $f$ and $F$ are positive, and $F(x)=\int f(\rho x)d\mu_R$.
The group $R_\gamma$ is unipotent and unimodular, so
$F(\rho x)=F(x)$.  Thus $\chi=f/F$ has the asserted property.
For every nonnegative measurable $R_\gamma$-invariant $F_0$,
the original Weil formula gives
\begin{equation}\label{eq:degenerate-cutoff-identity}
                      \int_G F_0\,d\nu_R
                          =\int_G\chi F_0\,dg.
\end{equation}
This applies to absolute scalar pairings; no phase compatibility
is inserted into the definition of the original operator.

\emph{The one-sided estimate and the compact dimension.}
For $R\ge1,y>0$ let
\[
 \mathcal E(R,y)=
  \{x:S_{z'}(x),S_z(\gamma x)\le R,\quad y<Y_\gamma(x)\le2y\},
\]
and put $Q_i=\{i=0\}\cup\{S_{z'}>2^iH\}$.
The first set is $R_\gamma$-invariant; the second is
$H_{z'}N$-invariant.  The latter fact uses that $H_{z'}$
fixes the original flag vectors, not that it is Euclidean
unitary in these coordinates.
For a unit vector $w\in\mathcal V'$, unfold
the left side of \eqref{eq:degenerate-cutoff-identity} with
$F_0={\bf1}_{\mathcal E(R,y)\cap Q_i}
|\langle w,\Phi'\rangle|^2$ over $H_{z'}N$.
The compact group preserves $dn$: its modulus on $N$ is
a continuous homomorphism to $\mathbb R_{>0}$ and is trivial.
Changing $hn$ to $nh$ and applying
Theorem~\ref{thm:original-degenerate-fibre-estimate} therefore gives
\begin{align*}
 &\int_{\mathcal E(R,y)\cap Q_i}
                  |\langle w,\Phi'(x)\rangle|^2d\nu_R(x)\\
 &\quad\le\frac{\kappa}{s}V(R,y)
       \int_{Q_i}\int_{H_{z'}}
                  |\langle w,\pi'(h)\Phi'(x)\rangle|^2dh\,d\nu_{z'}(x),
\end{align*}
where
\[
 V(R,y)=(4r)^{4r^2}R^{16r^3}\Delta^{8r^2}
                    \left(\frac{\pi\mathrm e}{a}\right)^{ja}(2y)^{2a}.
\]
The compact integral is exactly
$\|\Phi'(x)\|^2/\dim\mathcal V'$.
Indeed the average of the corresponding rank-one operators
commutes with the irreducible compact representation, so its
eigenspaces force it to be scalar; its trace determines that scalar.
The length tail, or the total mass when $i=0$, consequently gives
\begin{equation}\label{eq:actual-degenerate-one-side}
 \int_{\mathcal E(R,y)\cap Q_i}
       |\langle w,\Phi'\rangle|^2d\nu_R
 \le\frac{\kappa}{s}V(R,y)(4r+4)2^{-2qi}c_\Phi.
\end{equation}
The compact dimension has canceled exactly.
The fibre null-set assertion, followed by the same nonnegative
unfolding and cutoff, also shows $Y_\gamma>0$ $\nu_R$-almost
everywhere.  The reversed argument gives the corresponding
assertions for $\Phi$, $\gamma^{-1}$ and its transported measure.

\emph{The two projected Gram volumes.}
Write
\[
 S=S_z(\gamma x),\quad S'=S_{z'}(x),\quad
 Y=Y_{\gamma^{-1}}(\gamma x),\quad Y'=Y_\gamma(x).
\]
Exterior Cauchy--Schwarz for the first $l$ flag vectors gives
$|D_l(\gamma)|\le(SS')^l\le(SS')^r$ for $0\le l\le r$.
Projection decreases an exterior volume, so the two Gram
volumes also satisfy
\[
 Y,Y'\le\max(S,S')^r\Delta,\qquad
 \sigma\le\max\left(1,\frac{(SS')^r}{(2H)^{2r}}\right).
\]
In particular, the essential bound $D\le(SS')^r$ is retained.

We now prove the remaining product estimate
\begin{equation}\label{eq:actual-degenerate-Gram-product}
                                  YY'D\le2^j.
\end{equation}
In the original Euclidean space let
\[
 A=x^{-1}E_0,\quad B=x^{-1}\gamma^{-1}E_0,\quad
 C=x^{-1}\mathcal C_\gamma,\quad V_+=\ker(J-I).
\]
The two projection spaces are
$P=V_+\cap(C+A)$ and $Q=V_+\cap(C+B)$.
The form $h$ is positive semidefinite on $C$.
To see this, apply $x$: the common orthogonal space lies in
$E_0^{\perp_h}\cap z'^{\perp_h}$, whose radical is $E_0$
and whose middle part is the positive definite space
$z'^{\perp_{h_W}}\subset W$.

For $p\in P$, $p'\in Q$, choose decompositions
$p=c+a_0$, $p'=c'+b_0$ with $c,c'\in C$, $a_0\in A$,
$b_0\in B$.  The common space is orthogonal to both
isotropic spaces.  Since $p,p'\in V_+$,
\[
 h(c,c)=\|p\|^2,\qquad h(c',c')=\|p'\|^2,\qquad
 h(a_0,b_0)=\langle p,p'\rangle-h(c,c').
\]
The Cauchy--Schwarz inequality for the positive semidefinite
form on $C$ gives
\begin{equation}\label{eq:positive-common-space-pairing}
                         |h(a_0,b_0)|\le2\|p\|\|p'\|.
\end{equation}

Put $a_i=x^{-1}e_i$, $b_i=x^{-1}\gamma^{-1}e_i$ for
$1\le i\le j$, and write $Z_\gamma=(h(a_b,b_a))_{a,b\le j}$.
The common intersection $A\cap B$ has dimension $r-j$.
This follows first from the normal form and then by unipotent
transport.  The invertible cross matrix $Z_\gamma$ shows that
$a_1,\ldots,a_j$ and $b_1,\ldots,b_j$ complement that
intersection in $A$ and $B$.  The intersection is orthogonal
to $A+B$, so replacing $a_0,b_0$ by those complements does
not change their pairing.

Let $L_A:\mathbb C^j\to Q$ have columns $Q a_i$, and let
$L_B:\mathbb C^j\to P$ have columns $P b_i$, with $P,Q$
also denoting the orthogonal projections.  Solving for the
coordinates in the two complements gives exactly
\[
                  h(a_0,b_0)
                        =\langle L_A Z_\gamma^{-1}L_B^*p,p'\rangle.
\]
For example the first coordinate vector is
$Z_\gamma^{-1}(h(p,b_i))_i=Z_\gamma^{-1}L_B^*p$; the other
pairing row is $(h(a_i,p'))_i$.  Thus
\eqref{eq:positive-common-space-pairing} proves
$\|L_A Z_\gamma^{-1}L_B^*\|_{\rm op}\le2$.
Its $j$th exterior operator has norm
\[
 \frac{\|Qa_1\wedge\cdots\wedge Qa_j\|
          \|Pb_1\wedge\cdots\wedge Pb_j\|}{|\det Z_\gamma|}
                      =YY'D.
\]
This identity also holds if either family is dependent, when
both sides are zero.  The exterior operator norm is at most
$2^j$, proving \eqref{eq:actual-degenerate-Gram-product}.

\emph{The original dyadic estimate.}
For unit $v\in\mathcal V$, $w\in\mathcal V'$, set
$A_v(x)=|\langle v,\Phi(\gamma x)\rangle|$ and
$B_w(x)=|\langle w,\Phi'(x)\rangle|$.
In the dyadic estimate take
\[
 e=16r^3,\qquad m=a,\qquad
 K=K'=\frac{\kappa}{s}c_\Phi(4r+4)(4r)^{4r^2}\Delta^{8r^2}
                 \left(\frac{\pi\mathrm e}{a}\right)^{ja}4^a.
\]
For shell indices $i,i'$, use $R=2^{i+i'+1}H$ in
\eqref{eq:actual-degenerate-one-side} and in its reversed
counterpart.  This supplies exactly the two one-sided square
integrals in (OD3), with the specified $K,K'$.
All the pointwise requirements, including positivity almost
everywhere and $D\le(SS')^r$, were proved above.
Therefore
\[
 \int_G A_vB_w\,d\nu_R
 \le16K(2H)^{16r^3+4r}\Delta^2
       (2^j/D)^a\,\sigma^{-(q-16r^3-4r-1)/r}.
\]
The right side is finite.  Applying this to finite orthonormal
bases of $\mathcal V,\mathcal V'$ proves integrability of
$\|\Phi(\gamma x)\|\,\|\Phi'(x)\|$, hence of the original
rank-one kernel.  The bound is uniform for unit $v,w$;
dividing by $c_\Phi$ therefore bounds the operator norm,
without inserting either dimension.  Finally
$4^a2^{ja}=2^{(j+2)a}\le8^{ja}$ because $j\ge1$.
This gives \eqref{eq:actual-local-degenerate-bound}.
\end{proof}

\paragraph{The required moments fit the original arithmetic allowance.}
There is no further moment-range hypothesis in this local
application.  To state this precisely, retain the mathematical
standing conditions of
Theorem~\ref{global-gram-original-all-length-tail}, and denote
the original cloud parameter by $\delta_{\rm ar}$, to
distinguish it from the reciprocal-minor bound $\Delta$.
The standing data give
\[
 \delta_{\rm ar}\ge(kfL)^{C_\ell},\qquad
 h_{\rm NC}\ge C_\ell n^{19}f^5L,\qquad
 Q=\min\left(\frac{h_{\rm NC}}{10^{16}\sqrt m},
                                  \frac{4\delta_{\rm ar}}{10^8}\right).
\]
Every integer $1\le q\le C_1kr$ already satisfies both
extra numerical conditions of the length-tail theorem.
Indeed its proof gives $mq\le n^3$.
The interval hypothesis $B_v\ge10^8n^7$ gives
$B_v/(10^8n^4)\ge n^3$.
Since $m\le n$, $n\ge C_\ell\ge10^8$, $f,L\ge1$,
\[
 \frac{h_{\rm NC}}{10^{16}\sqrt m}
       \ge\frac{C_\ell n^{18}}{10^{16}}\ge n^3.
\]
Also $C_\ell\ge4$, $n\le2k$, and
$k\ge m^4\ge10^{32}$ give
\[
 \frac{4\delta_{\rm ar}}{10^8}
       \ge\frac{4k^4}{10^8}\ge8k^3\ge n^3.
\]
Thus $mq\le Q$ as well.  For $C_1\ge1$ the whole required
interval is nonempty: $k\ge(2r+2)^4$ and $r\ge1$ imply
\[
             16r^3+4r+2\le kr\le C_1kr.
\]
In particular $q=\lfloor C_1kr\rfloor$ is permitted.
The constants $c_q=10^{-8}$, $C_q=10^8$ and $C_H$ were
fixed before $C_1$, as in the original quantifier order.
We have enlarged neither a constructed tensor family nor
its shift range.  This verifies the actual moment input of
the local orbital theorem.

The remaining diagonal zero-Gram obligation is arithmetic:
assemble these local bounds with the actual leading-minor
count, stabilizer masses and global normalizers, and prove
absolute summability with the prescribed exponential saving.
The generic returns and final arithmetic Gram comparison
remain separate.
